% One-pager template — TUTORIAL / DOCUMENT digest layout.
% For teaching decks and reports: the goal is a map of the ideas, not a paper abstract.
% Keep to ONE page: build_onepager.py enforces it.
\documentclass[10pt]{article}
\usepackage[a4paper,margin=1.4cm]{geometry}
\usepackage{amsmath,amssymb,bm}
\usepackage{multicol}
\usepackage{enumitem}
\setlist{nosep,leftmargin=1.2em}
\usepackage{titlesec}
\titleformat{\section}{\bfseries\small}{}{0pt}{}
\titlespacing{\section}{0pt}{4pt}{2pt}
\pagestyle{empty}
\setlength{\parindent}{0pt}

\begin{document}
{\large\bfseries {{TITLE}}}\\[-2pt]
{\footnotesize {{SOURCE}} \hfill {{AUTHOR_OR_GROUP}} \hfill {{DATE}}}
\hrule\vspace{4pt}

\section*{In one line}
% The single takeaway. ~25 words.
{{ONE_LINE}}

\begin{multicols}{2}
\section*{Key ideas}
% The 3–6 concepts the material teaches, each in a clause.
{{KEY_IDEAS}}

\section*{Concept map}
% How the ideas depend on each other — prerequisites -> results.
{{CONCEPT_MAP}}

\section*{Worked takeaways}
% The equations/methods worth keeping, verbatim. Keep 2–4.
{{WORKED_TAKEAWAYS}}

\section*{Prerequisites}
% What you need to already know (from the learner profile: flag any not-yet).
{{PREREQUISITES}}

\section*{Where it fits}
% From context.md: the source's body of work / the field.
{{WHERE_IT_FITS}}

\section*{Relevance to my work}
% ~40 words: what to apply, linked notes.
{{RELEVANCE}}
\end{multicols}
\end{document}
